%% kapa: exam=1 round=33 year=2022 period=1 form=A subject=economics subject_ko=경제학원론 seq=26 number=66
%% source: Q-Net 감정평가사 기출문제 (cst00309) artlSeq=5209997 file=raw/1008/5209997/2022년도 제33회 감정평가사 1차시험 1교시(민법,경제학,부동산학).pdf
%% conversion: kordoc 4.24.51 -> tools/gen1.py (auto); review=proofread-sonnet; flags=none
\begin{question}{66}
아래의 개방경제 균형국민소득 결정모형에서 수출이 100만큼 늘어나는 경우 \uline{(ㄱ) 균형소득의 변화분}과 \uline{(ㄴ)경상수지의 변화분}은? (단, C는 소비, Y는 국민소득, T는 세금, I는 투자, G는 정부지출, X는 수출, M은 수입이며, 수출 증가 이전의 경제상태는 균형이다.)
\begin{framedtext}
○ C = 200 + 0.7(Y $-$ T) \hfill ○ I = 200 \hfill ○ G = 100 \par
○ T = 100 \hspace{4.2em} ○ X = 300 \hfill ○ M = 0.2(Y $-$ T) \hfill\mbox{}
\end{framedtext}
\begin{choices}
  \item ㄱ: 1000, ㄴ: 100
  \item ㄱ: 1000/3, ㄴ: 100/3
  \item ㄱ: 1000/3, ㄴ: 100
  \item ㄱ: 200, ㄴ: 60
  \item ㄱ: 200, ㄴ: 100
\end{choices}
\end{question}
\answer{4}
